\documentclass[10pt, conference]{IEEEtran}

% ---------- Reduce top margin / title block whitespace ----------
\usepackage[top=0.4in, bottom=0.4in, left=0.55in, right=0.55in]{geometry}

% ---------- Packages ----------
\usepackage{times}
\usepackage{titlesec}
\usepackage{enumitem}
\usepackage{xcolor}
\usepackage[colorlinks=true, linkcolor=black, citecolor=black, urlcolor=black]{hyperref}

% Force all body text to pure black
\color{black}

% ---------- Section formatting (compact, formal) ----------
\titlespacing*{\section}{0pt}{2pt}{1pt}
\titlespacing*{\subsection}{0pt}{1pt}{0pt}

\setlength{\parskip}{1pt}

% ---------- Title fields: EDIT PER PAPER ----------
\newcommand{\paperTitle}{Exploiting Page Table Locality for Agile TLB Prefetching}
\newcommand{\paperAuthors}{Vavouliotis, Alvarez, Karakostas, Nikas, Koziris, Jimenez, Casas}
\newcommand{\paperVenue}{ISCA 2021}

\begin{document}

\title{\Large\bfseries Paper Review: Exploiting Page Table\\Locality for Agile TLB Prefetching}

\author{
\IEEEauthorblockN{Vedang Kashyap}
\IEEEauthorblockA{
\textit{Department of Electronics Engineering (VLSI Design)}\\
\textit{Vellore Institute of Technology, Vellore, India}}
}

\maketitle

% ---------- 1. Introduction ----------
\section{Introduction}
Prior TLB prefetchers always trigger a background page walk to fetch a predicted PTE, so an inaccurate prefetch still costs a full memory-hierarchy traversal; none exploit that a page walk already returns a full cache line of adjacent PTEs for free, and no single prefetcher performs best across all workloads.

\subsection{Contribution}
\begin{itemize}[leftmargin=*, itemsep=2pt, topsep=2pt, parsep=1pt]
    \item Proposes Sampling-Based Free TLB Prefetching (SBFP), which uses sampling to predict which cache-line-adjacent PTEs fetched for free at the end of a page walk should be placed in the prefetch queue.
    \item Proposes Agile TLB Prefetcher (ATP), a composite prefetcher combining three low-cost prefetchers with logic that selects among them or disables prefetching per TLB miss.
\end{itemize}

% ---------- 2. Proposal ----------
\section{Proposal}
SBFP exploits that a page walk returns a 64-byte line holding the requested PTE plus 7 free neighbors, each at a distance of -7 to +7 excluding zero. A Free Distance Table of 14 saturating counters tracks each distance's hit rate: a free PTE enters the prefetch queue if its counter exceeds a threshold, else it goes to a Sampler tracking untrusted distances. Counters increment on a hit and decay on saturation, adapting as data structures change. ATP is a decision tree of three low-cost prefetchers, a stride, a two-distance, and a PC-indexed predictor, each with a Fake Prefetch Queue tracking its accuracy without real prefetches. Counters select which prefetcher fires, or none, based on these queues; ATP and SBFP share one prefetch queue.

% ---------- 3. Key Results ----------
\section{Key Results Achieved}
\begin{itemize}[leftmargin=*, itemsep=2pt, topsep=2pt, parsep=1pt]
    \item Geometric speedup over no prefetching is 16.2\%, 11.1\%, and 11.8\% for Qualcomm, SPEC, and GAP/XSBench workloads, respectively; over the best individual state-of-the-art prefetcher per suite, gains are 8.7\%, 3.4\%, and 4.2\%.
    \item Page walk memory references drop by 37\%, 26\%, and 5\% for the same three suites; total hardware cost is 1.68KB (ATP) plus 0.31KB (SBFP), with only 1.7-3.6\% of prefetches harmful to the OS page replacement policy.
\end{itemize}

% ---------- 4. Key Takeaway ----------
\section{Key Takeaway}
\begin{enumerate}[leftmargin=*, itemsep=2pt, topsep=2pt, parsep=1pt]
    \item A resource already paid for should be exploited before spending anything new, but only after measuring confidence in it rather than acting blindly; the 7 neighboring PTEs arrive free from a page walk that already happened, and the Free Distance Table and Fake Prefetch Queues exist purely to test which of them are worth using before committing real prefetch slots.
    \item No single fixed-feature predictor generalizes across workloads, so building one composite predictor that arbitrates between several narrow, cheap predictors can beat investing in one larger, more general predictor.
    \item Simulating a decision before committing to it is cheaper than committing and correcting; the Fake Prefetch Queues let every constituent prefetcher's accuracy be tracked concurrently without ever issuing a real, costly prefetch, turning selection into a side effect of ordinary operation rather than a separate probing phase.
\end{enumerate}

\section{Weakness}
\begin{itemize}[leftmargin=*, itemsep=2pt, topsep=2pt, parsep=1pt]
    \item Several design choices, the generalized Free Distance Table over a per-PC one, the fixed threshold, decay scheme, and counter widths, and ATP's three fixed constituent prefetchers, are all frozen at authoring time for complexity reasons, with no stated mechanism to re-tune or reconfigure them for workloads unlike those evaluated.
    \item For the GAP and XSBench workloads, ATP with SBFP reduces page walk memory references by only 5\%, far below the 37\% and 26\% seen for Qualcomm and SPEC, indicating the approach is markedly less effective on highly irregular big-data access patterns.
    \item ATP's throttling and selection logic is driven entirely by predicted-page accuracy in each Fake Prefetch Queue, with no signal for whether a correct prediction actually prevents a future TLB miss versus one that would have hit anyway, so selection optimizes accuracy rather than the metric that matters.
\end{itemize}

\section{Extension}
\begin{itemize}[leftmargin=*, itemsep=2pt, topsep=2pt, parsep=1pt]
    \item The paper explicitly leaves a per-data-structure Free Distance Table as unexplored beyond a single ideal-case test; a cheaper middle ground, such as a small set of tables selected by a coarse hash of the faulting instruction region rather than a full per-PC table, is untried.
    \item ATP's decision tree always fires exactly one constituent prefetcher or none; combining SBFP-style confidence counters with ATP's selection logic to allow two constituent prefetchers to fire together when both show high confidence, rather than forcing a single winner, is unexplored and could help on workloads where more than one access feature is present at once.
    \item Since the Fake Prefetch Queues already track each constituent prefetcher's predicted-page accuracy without real prefetches, the same infrastructure could log whether each fake hit corresponded to a page that would have missed the TLB anyway, giving ATP a truer usefulness signal than accuracy alone, at no added cost.
\end{itemize}

% ---------- 7. Reference ----------
\section{Reference}
Vavouliotis et al., \emph{Exploiting Page Table Locality for Agile TLB Prefetching}, ISCA 2021.

\end{document}
